\begin{table}[!htbp]
\centering
\small
\setlength{\tabcolsep}{3.2pt}
\begin{tabular}{>{\raggedright\arraybackslash}p{0.14\linewidth}>{\raggedright\arraybackslash}p{0.22\linewidth}>{\raggedright\arraybackslash}p{0.38\linewidth}>{\raggedright\arraybackslash}p{0.16\linewidth}}
\toprule
Choice & Audited contrast & Measured sensitivity & Reporting takeaway \\
\midrule
Tie rule & MMLU-Pro, OLMo-2 tokenizer
& \texttt{wrong}: 0.125914; \texttt{credit}: 0.532164 (40.6 pp, 4.2264$\times$).  The defensible \texttt{split}/\texttt{first}/\texttt{last} span is approximately 0.5--3.1 pp; \texttt{credit} exceeds intact \texttt{content\_norm} 0.207613 by 32.5 pp.
& Print the rule; treat \texttt{credit} only as an oracle bound. \\
Length unit & Characters versus continuation tokens
& \texttt{split} shifts by at most 2.02 pp; \texttt{credit} shifts by 14.53--35.20 pp.  OpenBookQA \texttt{credit} is 0.416 versus 0.644; ARC-Challenge tied-longest rates are 18.60\% versus 50.85\%.
& The convention alone is not reproducible specification. \\
Tokenizer & Item-matched token comparisons
& \texttt{split} shifts by 0.9003 pp; \texttt{credit} shifts by up to 10.6 pp on four-way tasks and 9.26 pp on MMLU-Pro.  The 32k/152k tokenizers yield fewer MMLU-Pro ties than the 100k/128k pair.
& Print the tokenizer; vocabulary size is not a monotone explanation. \\
Chance line & MMLU-Pro, $\texttt{n\_opt}\in\{3,\ldots,10\}$
& Naive $1/10=0.100000$ versus item-average 0.110877, against floor 0.116606: $+1.661$ pp and 1.1661$\times$ versus $+0.573$ pp and 1.0517$\times$.
& Print both defensible chance definitions. \\
\bottomrule
\end{tabular}
\caption{\textbf{Four under-specifications can move a content-side reference enough to reverse its reading.} The table moves the comparable numerical contrasts out of prose while preserving their scopes. Full per-convention floors are in Table~\ref{tab:conventions}; all values come from the same audited item records.}
\label{tab:reference-degrees}
\end{table}
